\section{The Gate}
\label{sec:gate}

The gate is the decision of Theorem~\ref{thm:gate}: after the draft, escalate
if and only if $L(\hat e_1-\hat e_2)>\lambda\,\Delta c$, where $\lambda$ is a
price of latency in real units that the user sets with one ``speed--care''
control (Proposition~\ref{prop:frontier}). It adds no model call to a request:
its features are by-products of the draft, and its own time is 0.08 to
0.12\,s per escalated request.

\subsection{Three rules, one decision record}
\label{sec:gate:rules}
\textbf{Verifier first.} On a verifiable task the exact verifier's verdict on
the draft decides: a passing draft is kept as \texttt{verified}, a failing one
is never kept. By Theorem~\ref{thm:precision} no confidence signal can add to
an exact verdict, so the cheap samples are skipped there.

\textbf{The fixed rule.} Without a verdict and without a promoted calibration,
the gate compares the draft's mean token log-probability with a threshold
derived from $\lambda$ ($-0.175$ at the default) and, inside an uncertainty
band of 0.10 around it, the agreement of $k=2$ cheap \Sone\ samples with the
draft. Unverified mathematics drafts always get their samples, whatever their
token confidence says, because their token confidence does not separate right
from wrong (Section~\ref{sec:eval:gate}). The samples are decoded in the
draft's own batch and dropped if the draft is stopped early.

\textbf{The calibrated rule.} A fitted calibration maps the features (mean,
minimum and lowest-window log-probability, agreement, domain, reask) to
$\hat e_1$; $\hat e_2$ is a per-domain estimate from the audit
stream. The calibration is fitted from the ledger's labelled drafts and is
\emph{promoted} to decide only when its held-out expected calibration error is
at most 0.05 on enough samples (Proposition~\ref{prop:miscal}); until then the
fixed rule decides and the calibrated estimate is only logged. A budget mode
solves for the multiplier $\lambda^\star$ that meets a target escalation rate
(Theorem~\ref{thm:gate}(ii)).

Every decision is recorded with the rule that made it, the features, the
estimates and the price, so a later refit can learn from it and a person can
read why a request was or was not escalated.

\subsection{Certified per-bin risk}
\label{sec:gate:bins}
A small calibration error bounds the \emph{average} miscalibration. A stronger
guarantee, taken from vLLM Semantic Router's cascade
\cite{semanticrouter2026}, bounds each score bin separately. The fitted
calibration keeps $B$ equal-mass bins of the held-out $\hat e_1$ (about one per
hundred drafts, two to five) with their counts and, for each, the one-sided
Clopper--Pearson \cite{clopper1934} upper bound on its error rate at level
$\alpha/B$: the smallest $u$ with $I_u(k{+}1,\,n{-}k)\ge 1-\alpha/B$, found by
bisection on the regularised incomplete beta function ($u=1-(\alpha/B)^{1/n}$
when $k=0$). A draft is kept only in a bin whose bound is at or below a
configured risk; otherwise the decision says \texttt{uncertified\_bin} with
the bound, and a calibration with no certified bin is not promoted. Our exact
verifiers supply the error counts at no cost. With $\alpha=0.05$ and five
bins, a bin of twenty drafts none of which was wrong still bounds its error at
0.21, so a young ledger certifies little and escalates more; the bound
tightens with the ledger by itself. The rule applies to the calibrated mode
only: budget mode fixes the escalation rate by design.

\subsection{Unknown features and repeated questions}
A draft without log-probabilities (a backend that does not return them, a
failed sample) must not look like a confident draft. The gate applies a
declared policy, \texttt{on\_unknown} = \texttt{escalate} (default),
\texttt{keep} or \texttt{fixed}, and the decision names the missing feature. A
user turn whose content words overlap the previous user turn's at Jaccard
$\ge0.6$ is a \emph{reask}: asking again is evidence that the last answer
failed. A reask of an unverified draft escalates under the fixed rule, and
the feature is logged so a refit can learn its weight. The test is lexical
rather than a cosine, so the gate runs without the embedding server.

\subsection{Sequential evidence}
\label{sec:gate:seq}
Three mechanisms read a stream of evidence and stop on a bound.

\begin{proposition}[Exact early escalation on a running minimum] \textnormal{[exact]}
\label{prop:n2}
Let $w(n)$ be the lowest mean log-probability over complete windows of $W$
consecutive tokens among the first $n$ tokens of a draft. Then $w$ is
non-increasing in $n$, so for the rule ``escalate iff $w(N)<\theta$'' the
stopping time $T=\min\{n:w(n)<\theta\}$ gives exactly the full draft's decision,
and no rule on the same stream escalates such a draft earlier without a
non-zero disagreement rate.
\end{proposition}
\begin{proof}
A minimum over a growing set; $w(N)\le w(T)<\theta$; before $T$ every window is
$\ge\theta$. \qed
\end{proof}
A draft that the gate will reject anyway can therefore be stopped at the
crossing and handed to \Stwo. On the mean-log-probability rule the same idea
is only a heuristic (a draft whose running mean is below the threshold minus a
margin is stopped), and Wald's bound shows it cannot be made more aggressive:
with the measured per-token spread $\sigma\approx0.28$ and margin 0.04, any
sequential test at 5\,\% errors needs about 260 tokens, against drafts of 47
to 291.

\begin{proposition}[Posterior stopping with a collapsing bound] \textnormal{[exact for the rule; strong for optimality]}
\label{prop:n4}
With $k$ probed answers clustered into categories, counts $v$ and a uniform
Dirichlet prior, the posterior probability that the plurality answer is the
majority answer is
$1-I_{1/2}(v_{\mathrm{top}}+1,\,k-v_{\mathrm{top}}+1)$.
Under a per-probe cost and a hard budget the Bayes-optimal rule stops at the
first $k$ with this probability at or above a bound $\theta(t_k)$ that is
non-increasing in the budget used \cite{frazier2007deadline}.
\end{proposition}

\begin{proposition}[Optimal probe interval] \textnormal{[exact under an exponential model]}
\label{prop:n5}
If the reasoning length at which the answer becomes available is exponential
with mean $m$ tokens, a probe costs $c_p$ seconds and decoding runs at $v$
tokens per second, the periodic interval minimising probe cost plus overshoot is
$\Delta^\star=\sqrt{2c_pvm}$.
\end{proposition}

\textbf{Adaptive reasoning exit.} The deliberate path is almost entirely
\Stwo\ reasoning, so the largest saving is to shorten the reasoning itself.
\Stwo\ reasons in chunks on a raw continuation endpoint, which keeps the
prefix cached. At a reflection cue (``Wait'', ``Alternatively''), with a
minimum gap of 64 tokens and at the latest every 256, the reasoning is closed
provisionally and the model is asked for its answer in a few tokens (0.6 to
0.8\,s). On a verifiable task the exact verifier judges that probe and, if it
passes, the reasoning stops; on factual tasks two consecutive agreeing probes
stop it, or optionally the posterior rule of Proposition~\ref{prop:n4}. The
final answer is then generated after the reasoning as on the standard path, so
the exit changes only how much reasoning is generated. The gap of 64 tokens is
within 1\,\% of the optimum of Proposition~\ref{prop:n5} at the measured
values ($c_p\approx0.7$\,s, $v\approx8$, $m\approx300$ give
$\Delta^\star\approx58$).

\subsection{Learning from its own traffic}
\label{sec:gate:labels}
The calibrated rule needs labelled drafts, and Proposition~\ref{prop:audit}
says routed traffic alone gives biased ones. Two mechanisms supply labels
without a \Stwo\ call. A \emph{probe} is a forced-fast draft that still draws
its agreement samples and is judged by the exact verifier. A \emph{label
stream} runs probes on verifiable practice tasks while no user request is in
flight; a user arriving cancels the draft in progress. At about 90 labels an
hour the 2,900 held-out samples that a five-bin calibration estimate needs
are about 32 hours of idle time. The quality page shows how many checked
answers are held and whether the learned rule has been promoted.
